\section{Synchronized operations on tensor products}\label{sec:lay}

A finite abelian group may have many cyclic factors of unrelated orders.
Applying a separate algorithm along each axis can charge for a full pass
through the data each time. We construct circuits that perform the same
kind of operation on many axes in one pass. This requires control of
both the coefficients and the workspace. A shear with entries in the
prime field acts separately on each base coordinate of an extension
scalar. A transform that restores its scratch registers can borrow them
from neighboring input arrays, even after those arrays have been
transformed. These two facts let us synchronize transforms over an
extension and charge their cost to the original data. They will be used
first for polynomial multiplication in Section~\ref{sec:pack} and then
for the Fourier upper bound in Section~\ref{sec:last}. For related
synchronization and generation arguments see
\cite[Sections~2--3]{FiniteFourier}.

\subsection{Shears and permutations}

Put \(U=\left(\begin{smallmatrix}1&1\\0&1\end{smallmatrix}\right)\).
It adds one coordinate to another. The iterated shear \(U^{\otimes j}\)
acts on \(2^j\) coordinates and gives the basic layer for combining
operations on different axes.

\begin{lemma}[Tensor powers of a shear]\label{lem:shear}
For each \(\beta>1/2\) and every \(j\ge0\), \(U^{\otimes j}\) has a
binary linear circuit of size \(O_\beta(2^j(1+j)^\beta)\) over every
field, with an implied constant independent of the field.
\end{lemma}

\begin{proof}
Work first over an algebraic closure of the prime field. Use radix
\(\ell=2\) except in characteristic two, where we use \(\ell=3\).
On four coordinates consider the group generated by invertible
monomial matrices (permutations and diagonal scalings) and the Fourier
matrix \(D_\ell\) on any \(\ell\) of the coordinates, fixing the rest.
This group is \(\mathrm{GL}_4\). To see this, consider the tangent
spaces at the identity of its conjugates of the diagonal subgroup.
They contain all diagonal directions. Given \(i\ne j\), choose a
Fourier block containing \(i,j\). Conjugating a diagonal matrix unit
in the block by its Fourier matrix gives a matrix with a nonzero
\(ij\)-entry, because all entries of the transform and its inverse are
nonzero. Further conjugation by \(\operatorname{diag}(x_1,\ldots,x_4)\)
multiplies the \(ij\)-entry by \(x_i/x_j\). The distinct off-diagonal
characters and the constant character are linearly independent over
the algebraic closure. Taking linear combinations of these conjugates therefore isolates
the matrix unit in each off-diagonal position.

Choose finitely many of these tori whose tangent spaces span the space of
\(4\times4\) matrices. The multiplication map from their product to
\(\mathrm{GL}_4\) has surjective differential at the identity tuple.
Its image contains a nonempty open subset. Any translate of this open
subset meets it by irreducibility, so the generated subgroup equals
\(\mathrm{GL}_4\). In particular \(U^{\otimes2}\) has a finite word
in the generators. Inverses are again words, since Fourier inversion
uses only a permutation and a diagonal scaling in addition to the
forward transform.

The word length and the degree of its coefficient field over the prime
field can be bounded uniformly. Choose one word in characteristic zero.
Its coefficients lie in a number field and satisfy the matrix identity
for \(U^{\otimes2}\). Outside finitely many rational primes we may
reduce all coefficients at a prime of that number field, preserving
the identity and all nonzero scalars that are inverted. The residue
degrees are bounded by the number-field degree. In each exceptional
characteristic use the construction above.

For \(j=2r\), take the \(r\)th tensor power of this word over the
compositum with the given base field. Each monomial letter remains
monomial. A Fourier letter acts as \(D_\ell\) on a block and as the
identity on its complement. Its tensor power is a direct sum of
blocks \(D_\ell^{\otimes k}\), \(0\le k\le r\), obtained by fixing
the positions and labels outside the Fourier block. Their dimensions
sum to \(4^r\). Theorem~\ref{thm:kernel} therefore bounds the whole
letter by \(O_\beta(4^r(1+r)^\beta)\) gates. Describing the bounded-degree
extension in a base-field basis simulates every gate at constant cost.
The output matrix has base-field entries. For odd \(j\) apply the
remaining \(U\) on pairs of coordinates. The identity handles \(j=0\).
\end{proof}

A \emph{word} will be a sequence of invertible layers over an extension
field \(E\). A monomial layer permutes and rescales coordinates. A
\emph{shear layer} partitions the coordinates into blocks and applies
\(U^{\otimes j}\) to each block of size \(2^j\), allowing a different
\(j\) for each block. Singleton blocks are fixed. Permutations may
arrange the coordinates into the required blocks.

Tensoring two monomial layers gives a monomial layer. Tensoring two
shear blocks gives \(U^{\otimes(j_1+j_2)}\), so tensoring shear layers
gives a shear layer up to coordinate permutations. To tensor words of
length at most \(D\), assign a monomial slot and a shear slot to each
step and put the identity into each unused slot. The synchronized word
has length \(O(D)\), regardless of how many words are tensored. The
same schedule works for disjoint calls. An inverse word also has length
\(O(D)\), since \(U^{-1}\) is diagonally conjugate to \(U\).

When \(E\) contains a primitive \(s\)th root for \(s=\ell^j\), radix
recursion gives \(F_s\) a word of length \(O(1+\log s)\) on its \(s\)
coordinates \cite{CooleyTukey}. Each stage consists of disjoint
\(\ell\)-coordinate Fourier operations, diagonal scalings and
permutations. Gaussian elimination expresses any invertible operation
of this fixed size by a bounded number of monomial matrices and
elementary shears. A nonzero scalar shear is diagonally conjugate to
a unit shear.

\subsection{Evaluation and interpolation on a geometric progression}

For \(q\in K^\times\), the matrix
\(V_a(q)=(q^{ij})_{0\le i,j<a}\) evaluates a polynomial in the
power basis at \(1,q,\ldots,q^{a-1}\). The Newton basis converts both
evaluation and interpolation into a bounded number of triangular
Toeplitz operations over \(K\).

\begin{lemma}[Geometric Newton factorization]\label{up:newton-lemma}
Let \(V=V_a(q)\) have distinct nodes over \(K\). There is a lower
triangular matrix \(L\), diagonally equivalent over \(K\) to an
invertible triangular Toeplitz matrix, such that with
\(\Delta=\operatorname{diag}(L_{ii})\),
\begin{equation}\label{up:newton}
                         V=L\Delta^{-1}L^{\mathsf T}.
\end{equation}
Thus \(V\) and \(V^{-1}\) are bounded products of invertible diagonals
and invertible triangular Toeplitz matrices.
\end{lemma}

\begin{proof}
The \(j\)th Newton polynomial is \(\prod_{s<j}(X-q^s)\), where
\(0\le s<j<a\). Collect its power-basis coefficients in column \(j\)
of a unit upper triangular matrix \(P\). Its values form column \(j\)
of \(L=VP\), and vanish above the diagonal. Write
\(H_i=\prod_{k=1}^i(1-q^k)\), with \(H_0=1\). Evaluating the product
gives
\[
       L_{ij}=(-1)^j q^{\binom j2}H_i/H_{i-j}
               \qquad (0\le j\le i<a).
\]
Distinctness of the nodes makes each \(H_i\) nonzero. The lower
triangular Toeplitz matrix with coefficients \(1/H_k\) is therefore
diagonally equivalent to \(L\).

Now \(V=(L\Delta^{-1})\Delta P^{-1}\) is a unit
lower--diagonal--unit upper decomposition. Transposing and using
uniqueness for the symmetric matrix \(V\) gives
\(P^{-1}=(L\Delta^{-1})^{\mathsf T}\), proving \eqref{up:newton}.
An invertible lower triangular Toeplitz matrix represents multiplication
by a unit in \(K[X]/(X^a)\), whose reciprocal represents the inverse.
Transposition gives the upper triangular case. For this Newton change
of basis see \cite[Section~5]{BostanSchost}, for Vandermonde
factorizations see \cite{orucPhillips2000,kuznetsov2018}, and for their
implementation in circuits see \cite{FiniteFourier,ExplicitFourier}.
\end{proof}

\subsection{Using and restoring occupied workspace}

An axis transform will need scratch registers even when every coordinate
already holds data. The condition \(T\oplus I\) requires a word to
send \((x,z)\) to \((Tx,z)\) for \emph{every} initial value of the
scratch array \(z\). The identity on \(z\) makes a neighboring fiber
available as temporary storage. For earlier uses of restoration see
\cite{bennett1973,BuhrmanCleveKouckyLoffSpeelman2014,ExplicitFourier}.

\begin{lemma}[One factor with occupied workspace]\label{up:toeplitz}
Suppose \(2a\le b\le6a\), \(E\) contains a primitive \(b\)th root \(u\),
and a word of length \(D\) acts as \(F_b\oplus I\) on at most \(Cb\)
registers, where \(F_b=V_b(u)\). Let \(T\) be an invertible triangular
Toeplitz matrix of side \(a\), a geometric evaluation matrix with
distinct nodes, or the inverse of such an evaluation matrix, over
\(E\). A word of length \(O(D+1)\) acts as \(T\oplus I\) on
\(O((C+1)a)\) registers.
\end{lemma}

\begin{proof}
First implement \((x,y)\mapsto(x,y+Ax)\) for an arbitrary Toeplitz
matrix \(A\) of side \(a\). Write \(A=RLJ\), where \(J\) pads to \(b\)
coordinates, \(R\) truncates to \(a\), and \(L\) is a circulant.
The \(2a-1\) Toeplitz diagonals occupy distinct residues modulo \(b\).
For \(L\), transform the source and destination by \(F_b\), add the
scaled source entries using the diagonal \(F_b L F_b^{-1}\), and undo
the two transforms. The additions take one layer of scalar shears,
implemented by diagonal conjugation and unit shears. For \(J\) and
\(R\) the corresponding additions use disjoint unit shears.

For any two composable maps \(A_1,A_2\), the updates
\[
 z\gets z+A_1x,\qquad y\gets y+A_2z,\qquad
 z\gets z-A_1x,\qquad y\gets y-A_2z
\]
restore \(z\) and add \(A_2A_1x\) to \(y\). The contributions from the
original value of \(z\) cancel. Applying this identity twice implements
the addition for \(RLJ\), with a bounded number of length-\(b\)
transforms and \(O((C+1)b)\) registers.

To replace \(x\) by \(Tx\), first take \(T\) lower triangular Toeplitz.
Its action is multiplication by a unit \(h\) in \(E[X]/(X^a)\).
After dividing out the main diagonal we have \(h(0)=1\). Solve
\(f(uX)=h(X)f(X)\pmod{X^a}\) with \(f(0)=1\). The coefficient \(f_i\)
is determined by the earlier coefficients since \(u^i-1\ne0\) for
\(1\le i<a\). Write \(S\) for multiplication by \(f\) and
\(D_u=\operatorname{diag}(1,u,\ldots,u^{a-1})\). The identity for \(f\)
expresses \(T\) as \(D_u S D_u^{-1}S^{-1}\).

The additions associated to \(S\) and \(S^{-1}\) are available since
both matrices are Toeplitz. On two arrays they realize
\[
 \begin{pmatrix}I&S\\0&I\end{pmatrix}
 \begin{pmatrix}I&0\\-S^{-1}&I\end{pmatrix}
 \begin{pmatrix}I&S\\0&I\end{pmatrix}
 \begin{pmatrix}0&-I\\I&0\end{pmatrix}
  = B,\qquad B=\operatorname{diag}(S,S^{-1}).
\]
The last factor is a signed interchange. In
\((D_u\oplus I)B(D_u^{-1}\oplus I)B^{-1}\) the second block returns to
its initial value and the first block undergoes \(T\). Restore the main
diagonal by scaling. Reversing coordinates handles upper triangular
factors. For geometric factors and their inverses, successively apply
the words for the factors of Lemma~\ref{up:newton-lemma} using the same
scratch registers.
\end{proof}

\subsection{Borrowing registers across fibers}

Tensoring words that use \(Ca_i\) registers on the \(i\)th axis would
allocate as many as \(C^r\prod_{i=1}^r a_i\) registers. Instead we
borrow storage inside blocks of a fixed number of axes. A fiber along
one axis consists of the entries with every other coordinate fixed.

\begin{lemma}[Synchronizing axis operations]\label{up:fibers}
For each \(i\), suppose \(T_i\) acts on \(a_i\ge2\) coordinates and
a word of length at most \(D\) acts as \(T_i\oplus I\) on at most
\(Ca_i\) registers. Put \(B=\prod_i a_i\). A word of length
\(O_C(D+1)\) acts as \((\bigotimes_i T_i)\oplus I\) on \(O_C(B)\)
registers.
\end{lemma}

\begin{proof}
Put \(m=\lceil C\rceil\) and \(k=1+\lceil\log_2 m\rceil\). With at
least \(k\) factors, form blocks of \(k\) to \(2k-1\) factors. Within
one block, an axis of length \(a\) has at least \(2^{k-1}\ge m\)
fibers. Partition them into sets of \(m\) to \(2m\) fibers. Apply the
axis word to one fiber in each set, drawing its scratch registers from
the other fibers of that set. When the operation ends all borrowed
entries are restored. It is therefore possible to process the next
fiber even when the scratch space contains an earlier fiber's output.

Process all sets in parallel and the \(O_C(1)\) fibers per set in
sequence. Repeat for the \(O_C(1)\) axes of the block. This gives a
word of length \(O_C(D+1)\) on exactly the coordinates of each block.
Synchronizing their tensor product gives the same length bound on the
\(B\) data coordinates. With fewer than \(k\) factors, simply transform
one axis at a time, drawing disjoint scratch arrays from \(O_C(B)\)
additional registers for the simultaneous fiber calls. Every call
restores its array.
\end{proof}
